% Source: https://artofproblemsolving.com/community/c3h1682125_geometry
% Stable identity: aops-topic-1682125-q01
The radius of the second circle is $PA=PB=5$. In triangle $OPA$ the sides are $12,5,13$ and $12^2+5^2=13^2$, so $\angle OPA=90^\circ$.

\begin{asy}
usepackage("amsmath");
size(7cm);
defaultpen(fontsize(10pt));
pair O=(0,0),P=(12,0),A=(12,5),B=(12,-5);
draw(circle(O,13),linewidth(0.8));draw(circle(P,5),linewidth(0.8)+blue);
draw(O--P,linewidth(0.7));draw(O--A,dashed);draw(A--B,red+linewidth(0.9));
real r=0.9;draw(P+(-r,0)--P+(-r,r)--P+(0,r),linewidth(0.6));
dot(O);dot(P);dot(A);dot(B);
label("$O$",O,W);label("$P$",P,dir(-135));label("$A$",A,NE);label("$B$",B,SE);
label("$12$",(O+P)/2,S);label("$13$",(O+A)/2,NW);label("$5$",(P+A)/2,E);
label("\mbox{}",(min(currentpicture,user=true)+max(currentpicture,user=true))/2,black); // keeps the text after the figure black when LaTeX typesets the labels inline
\end{asy}

The common chord $AB$ is perpendicular to the line of centres $OP$ (both $O$ and $P$ are equidistant from $A$ and $B$, so $OP$ is the perpendicular bisector of $AB$). Since $PA\perp OP$, the chord $AB$ lies along line $PA$ and passes through $P$, and $P$ is its midpoint. Hence $AB=2PA=\boxed{10}$.
